\section{Chapitre~\ref{sec:neurontypes}. Types de neurones}

Les nombres du chapitre reposent sur des comptages directs de cellules dans le cerveau humain, là où ils existent; la règle \og un neurone, un neurotransmetteur\fg, sur des atlas cellulaires et sur des expériences qui en ont aussi trouvé les exceptions; le rôle de \nt{DOPA}, sur des enregistrements d'impulsions; et les rôles des autres canaux de diffusion, sur des hypothèses.

{\footnotesize
\setlength{\tabcolsep}{3pt}\renewcommand{\arraystretch}{1.15}
\begin{longtable}{>{\raggedright}p{0.5cm}>{\raggedright}p{4.0cm}>{\raggedright}p{5.6cm}>{\raggedright}p{2.3cm}>{\raggedright\arraybackslash}p{1.7cm}}
\toprule
N\textsuperscript{o} & Affirmation & Expérience et ce qui est mesuré & Source & Statut \\
\midrule\endfirsthead
\toprule
N\textsuperscript{o} & Affirmation & Expérience et ce qui est mesuré & Source & Statut \\
\midrule\endhead
1 & Le cerveau humain compte environ $8.6\cdot10^{10}$ neurones, et presque tous les neurones \nt{GLUT} sont les petits neurones du cervelet (tabl.~\ref{tab:types}) & Fractionneur isotrope: le tissu est dissous en une suspension homogène de noyaux, et l'on compte les noyaux porteurs du marqueur neuronal NeuN. Chez l'homme adulte, $86.1\pm8.1$ milliards de neurones; le cortex cérébral n'en contient que $19\,\%$, la masse principale est dans le cervelet & \cite{Azevedo2009} & mesuré \\
2 & Presque chaque neurone a un seul neurotransmetteur principal (proposition~\ref{prop:dale}), mais il y a des exceptions & Atlas transcriptomique du cerveau de la souris: environ $4\cdot10^6$ cellules, $5322$ groupes; chaque groupe reçoit un neurotransmetteur d'après les gènes de synthèse et d'empaquetage. Les exceptions sont mesurées directement: en tranches, les neurones \nt{DOPA} libèrent aussi du GABA par le même transporteur vésiculaire et inhibent les neurones du striatum; dans une partie des axones \nt{DOPA}, le glutamate et la dopamine sortent de terminaisons différentes d'un même fil (microscopie électronique et optogénétique) & \cite{Strata1999,Yao2023,Tritsch2012,Zhang2015} & mesuré \\
3 & Toute la colonne de la matrice des poids est de même signe~\eqref{eq:dalesign} & Déduction du chapitre à partir de la proposition~\ref{prop:dale} et du fait que presque tous les récepteurs du glutamate sont excitateurs et ceux du GABA inhibiteurs. Il n'existe pas de vérification directe de \og tous les destinataires d'un neurone reçoivent un poids de même signe\fg{} comme règle générale; contre-exemples: la co-libération de GABA par les neurones \nt{DOPA} (ligne~2) et les destinataires dotés de récepteurs de signes différents (ligne~11) & déduction du chapitre & théorie \\
4 & De trois quarts à quatre cinquièmes des neurones du cortex sont \nt{GLUT}, d'un cinquième à un quart sont \nt{GABA} & Immunomarquage du GABA dans des colonnes de $50$\,\textmu m de large traversant toute l'épaisseur du cortex, dans $10$ aires corticales du singe: dans $8$ aires, les neurones GABA forment environ $25\,\%$ de tous les neurones, dans le cortex visuel primaire moins de $20\,\%$ & \cite{Hendry1987} & mesuré \\
5 & Un neurone \nt{GLUT} a environ $10^4$ sorties & Stéréologie post mortem: dans le néocortex de jeunes hommes, $1.64\cdot10^{14}$ synapses (microscopie électronique, dissecteur) et environ $2.3\cdot10^{10}$ neurones. Le rapport donne $\sim7\cdot10^3$ synapses par neurone; en moyenne, le nombre d'entrées égale le nombre de sorties & \cite{Tang2001synapses,Pakkenberg1997} & mesuré \\
6 & La sortie d'un neurone \nt{DOPA} se ramifie en $10^5$--$10^6$ terminaisons; c'est une estimation d'après la couverture de la cible, pas un comptage & Marquage viral de la membrane de neurones \nt{DOPA} isolés du rat et reconstruction complète de l'axone: un neurone couvre de $0.45$ à $5.7\,\%$ (en moyenne $2.7\,\%$) du volume du striatum. C'est la couverture qui est mesurée, pas le nombre de terminaisons: celui-ci en est déduit en ordre de grandeur, sans comptage direct. Pour \nt{SERO} et \nt{NORA}, les nombres de terminaisons sont des estimations grossières & \cite{Matsuda2009} & mesuré \\
7 & Les neurones de diffusion sont peu nombreux: \nt{DOPA} $\sim5\cdot10^5$ (le principal des deux groupes, borne inférieure), \nt{SERO} $\sim3\cdot10^5$ (somme de deux comptages partiels), \nt{HIST} $\sim6\cdot10^4$ (tabl.~\ref{tab:types}, fig.~\ref{fig:typepie}) & Comptages stéréologiques chez l'homme: $550\,000$ neurones pigmentés dans la substance noire (sans l'aire tegmentale ventrale); $235\,000$ neurones dans le noyau dorsal du raphé (tous les neurones du noyau) et encore $125\,000$ neurones sérotoninergiques dans le tegmentum du pont; environ $64\,000$ neurones histaminergiques de l'hypothalamus. Pour \nt{NORA}, \nt{GLYC}, \nt{ACET} et \nt{ADRE}, pas de comptage direct: le tableau donne des estimations grossières en ordre de grandeur & \cite{Pakkenberg1991,Baker1990,Baker1991,Airaksinen1991} & mesuré \\
8 & Le glutamate dans la fente monte jusqu'à environ une millimole par litre et retombe en $\sim1\ms$ & D'après le déplacement d'un bloqueur compétitif à dissociation rapide des récepteurs NMDA pendant la transmission synaptique (culture de neurones d'hippocampe): pic de $1.1$\,mM, constante de décroissance de $1.2\ms$ & \cite{Clements1992} & mesuré \\
9 & Dans le cortex en activité, les courants excitateur et inhibiteur sont presque équilibrés, et le neurone reste près du seuil & Théorie: un réseau à connexions fortes atteint de lui-même le régime équilibré. Expérience: dans le cortex préfrontal du furet in vivo, pendant les états \og up\fg, le potentiel d'inversion du courant synaptique reste vers $-37$\,mV pendant des centaines de millisecondes, alors que la conductance varie de $\sim21$\,nS: excitation et inhibition montent et descendent ensemble & \cite{vanVreeswijk1996,Haider2006} & mesuré \\
10 & Les neurones \nt{DOPA} annoncent l'erreur de prédiction de la récompense~\eqref{eq:rpe} (fig.~\ref{fig:rpe}) & Impulsions des neurones \nt{DOPA} du singe: bouffée au jus inattendu; après apprentissage, bouffée au signal et aucune réponse au jus prédit; creux de fréquence quand le jus promis n'est pas donné. Dans l'expérience quantitative, la fréquence est proportionnelle à la différence entre la récompense actuelle et la moyenne pondérée des précédentes, mais seulement pour les erreurs positives. L'équation~\eqref{eq:rpe} elle-même est l'algorithme des différences temporelles & \cite{Schultz1997,Bayer2005,SuttonBarto2018} & mesuré \\
11 & Le signe et la vitesse sont fixés par le récepteur: ionotropes, fractions de milliseconde; métabotropes, beaucoup plus lents (proposition~\ref{prop:receptor}) & Brèves impulsions de glutamate sur des fragments de membrane de neurones d'hippocampe: le courant des récepteurs ionotropes monte (de $20$ à $80\,\%$) en $0.2$--$0.6\ms$ et retombe en $2$--$3\ms$. Le potentiel inhibiteur lent des neurones pyramidaux est supprimé par un bloqueur sélectif du récepteur métabotrope du GABA. Deux familles de récepteurs de la dopamine à action opposée; dans le striatum, les neurones à récepteur D1 ont besoin de dopamine pour renforcer les synapses, ceux à D2 pour les affaiblir & \cite{Colquhoun1992,DutarNicoll1988,Beaulieu2011,Shen2008} & mesuré \\
12 & Un canal de diffusion n'est pas un fil unique, mais un faisceau de quelques lignes (section~\ref{sec:buses}) & Marquage viro-génétique des neurones sérotoninergiques du noyau dorsal du raphé de la souris: les neurones qui projettent vers l'amygdale et vers le cortex frontal occupent des places différentes du noyau, se ramifient vers des régions complémentaires et répondent de façon opposée à un stimulus désagréable. Revue: des sous-groupes de neurones du locus coeruleus (\nt{NORA}) codent des processus différents & \cite{Ren2018,Poe2020} & mesuré \\
13 & \nt{NORA} fixe le gain $g$ & Hypothèse du gain adaptatif; modèle de réseau dans lequel les catécholamines augmentent la pente de la caractéristique de transfert. Pas de mesure directe de \og $g$ croît avec la noradrénaline\fg{} à l'échelle du réseau entier; mesuré: le diamètre de la pupille suit les impulsions des neurones du locus coeruleus du singe & \cite{AstonJones2005,ServanSchreiber1990,Joshi2016} & hypothèse \\
14 & \nt{ACET} commute \og écriture/lecture\fg{} et fixe la vitesse d'apprentissage & Hypothèse. Appui: dans des tranches de cortex olfactif du rat, les agonistes de l'acétylcholine ($100$\,\textmu M) affaiblissent les synapses internes \og de rappel\fg{} de plus de $60\,\%$, et les synapses d'entrée de moins de $15\,\%$ & \cite{Hasselmo2006,HasselmoBower1992} & hypothèse \\
15 & \nt{SERO} fixe le facteur d'actualisation $\gamma$; les neurones sérotoninergiques travaillent régulièrement, à quelques impulsions par seconde, et se taisent presque pendant l'une des phases du sommeil & Hypothèse \og sérotonine $=\gamma$\fg. Argument le plus fort: l'activation optogénétique des neurones sérotoninergiques du noyau dorsal du raphé de la souris réduit le nombre d'abandons prématurés de l'attente d'une récompense différée. La fréquence des impulsions et le silence pendant le sommeil à mouvements oculaires rapides sont mesurés (revue d'enregistrements) & \cite{Doya2002,Miyazaki2014,JacobsAzmitia1992} & hypothèse \\
\bottomrule
\end{longtable}}
